\documentclass{situationist-preprint}

\title{Intentions Are Committed Plans}
\msid{ACT-06-a}
\msfamily{ACT-06}
\msdiscipline{Action, Agency and Free Will}
\mstype{Principal manuscript}
\mskeywords{intention; planning theory; commitment; belief--desire model; cognitivism; Buridan cases; means--end coherence}

\begin{document}
\maketitle

\begin{abstract}
What is it to intend to do something? Reductive views identify intention with a combination of more basic attitudes: a predominant desire together with a belief about what one will do, or a belief that one will act formed in a distinctive way. We argue that intentions are plans held with commitment, and that they do not reduce to beliefs and desires, which involve no commitment. Commitment has two components: an intention settles a practical question, so that the agent is disposed not to reopen it without new reason, and it constrains further reasoning, by excluding incompatible options from consideration and by requiring the agent to settle on means. We prove two non-supervenience results. In a Buridan case, two agents with the same beliefs and desires intend different things, so intention is not fixed by belief and desire; and an agent who most wants to do something and expects to do it may still not have decided, so belief and desire do not suffice for intention. We then show that the means--end requirement governs intention but not the corresponding beliefs and desires, and answer the cognitivist reply. The account fits the restriction of agglomeration to single plans in ACT-15 and the belief constraint of ACT-11.
\end{abstract}

\section{The problem}

The belief--desire model of practical reason, in the form given to it by \citet{Davidson1963}, explains action by a pair of attitudes: a pro-attitude toward an end and a belief that an action is a means to it. When Davidson turned to intention, he first treated \enquote{with the intention of} as a mere relational expression, adding nothing to the belief--desire explanation, and later recognized pure intending, intending with no action yet under way, as a distinct state, which he identified with an all-out evaluative judgment that the action is desirable \citep{Davidson1978}. Others identified intention with a predominant desire together with a belief that one will act \citep{Audi1973}, or defended the belief--desire model as sufficient on its own \citep{Sinhababu2013}. Cognitivists identified intention with a belief that one will act, formed in a certain way \citep{Harman1976,Velleman1989,Setiya2007}.

\citet{Bratman1987} argued that intentions are elements of partial plans, and that their characteristic role in practical reasoning, coordinating the agent's conduct over time and with others, requires that they be distinct from desires and beliefs. We defend Bratman's thesis and give it a sharper argument. The argument has two parts: a pair of non-supervenience results, and a normative difference. Kavka's toxin puzzle, which concerns whether one can intend an action one has reason to intend but no reason to perform \citep{Kavka1983}, is taken up in ACT-06-b.

\section{Framework}

\begin{definition}[Plan]\label{d:plan}
A \emph{plan} is a structured content: an end together with means to it, sub-plans for those means, and times, typically partial in that some means and details are left to be filled in later.
\end{definition}

\begin{definition}[Commitment]\label{d:commit}
An attitude toward a plan involves \emph{commitment} if it has two features. (i) \emph{Settling}: the agent treats the question whether to execute the plan as closed, and is disposed not to reconsider it in the absence of new information or of a change in her reasons. (ii) \emph{Constraint}: the attitude governs further practical reasoning by excluding from consideration options incompatible with the plan, and by requiring the agent to settle on means to its ends.
\end{definition}

\begin{definition}[Intention]\label{d:intention}
An \emph{intention} is an attitude whose content is a plan and which involves commitment.
\end{definition}

Definition~\ref{d:intention} is the thesis, and the question is whether some combination of beliefs and desires could play the same role. To make the question precise:

\begin{definition}[Reduction]\label{d:reduction}
Intention \emph{reduces} to belief and desire if there is a condition $R$, stated only in terms of an agent's beliefs, desires and their strengths, such that necessarily an agent intends to $A$ if and only if she satisfies $R$ with respect to $A$.
\end{definition}

\begin{principle}[Functional individuation]\label{p:func}
Attitudes are individuated by their functional and normative roles: two attitude types are distinct if one is governed by a rational requirement that does not govern the other.
\end{principle}

\begin{assumption}[No commitment in belief and desire]\label{a:nocommit}
Neither a desire nor a belief, nor a combination of them, involves commitment. One may rationally desire incompatible things, desire something without settling whether to pursue it, and reopen the question how much one wants something whenever one likes; and a belief about what one will do does not, as such, exclude options from practical consideration or require means.
\end{assumption}

Assumption~\ref{a:nocommit} is the claim that the reductive views deny, and the argument of Section~\ref{s:arg} is designed to support it from cases rather than assume it.

\section{The argument}\label{s:arg}

\begin{proposition}[Buridan]\label{p:buridan}
Intention does not supervene on belief and desire: there are possible agents with the same beliefs and desires, of the same strengths, one of whom intends to $A$ and the other of whom intends to $B$.
\end{proposition}

\begin{proof}
A walker comes to a fork, either branch of which leads home in the same time by equally pleasant routes. She wants to get home, wants each route to the same degree, and believes that each will get her there. Being sensible, she does not stand at the fork: she settles on the left branch. A second walker with exactly the same beliefs and desires settles on the right. Each has an intention the other lacks, and since their beliefs and desires are the same, no condition $R$ stated in terms of beliefs and desires is satisfied by one with respect to the left branch and not by the other.
\end{proof}

The proof uses the possibility of picking among equally good options \citep{UllmannMargalitMorgenbesser1977}. It is no objection that, once the walker has settled, she comes to believe that she will take the left branch and may even come to want it more. These are effects of the intention, and the question is what it is to intend before they arise. The reductionist may reply that the walker's intention just is the belief she forms at the moment of settling, which is the cognitivist proposal answered below.

\begin{proposition}[Undecided expectation]\label{p:undecided}
Predominant desire together with the belief that one will act does not suffice for intention.
\end{proposition}

\begin{proof}
An agent is invited to a party next week. She wants to go more than she wants anything incompatible with going, and she expects, on the basis of her knowledge of herself, that she will go. But she has not decided: she is waiting to hear whether a friend will be there, and she treats the question as open. She satisfies the predominant-desire and belief conditions without settling, so without intending. If she were asked whether she intends to go, the correct answer would be \enquote{I haven't decided yet, though I expect I will go}, and the answer is not incoherent.
\end{proof}

\begin{theorem}\label{t:main}
Intention does not reduce to belief and desire.
\end{theorem}

\begin{proof}
If intention reduced to belief and desire, there would be a condition $R$ on beliefs, desires and their strengths that held of exactly the agents who intend. By Proposition~\ref{p:buridan}, two agents may agree on every such condition and differ in what they intend, so no $R$ is necessary and sufficient for intending to take the left branch. Proposition~\ref{p:undecided} shows, independently, that the most natural candidate for $R$ is not sufficient.
\end{proof}

The theorem would leave open that intention differs from belief and desire only in some causal respect without normative significance. It does not. Intentions are governed by a requirement that does not govern the beliefs and desires associated with them.

\begin{proposition}[Means--end requirement]\label{p:means}
An agent who intends to $E$, believes that $M$ is necessary for $E$, and does not intend to $M$, is incoherent, unless she revises the intention or the belief. An agent who desires $E$, even predominantly, believes that $M$ is necessary for $E$, and does not desire $M$ or intend to $M$, is not thereby incoherent.
\end{proposition}

\begin{proof}
For the first claim: intending $E$ settles the question whether to pursue $E$ (Definition~\ref{d:commit}(i)), and an end settled on while the necessary means remain unsettled is an unfinished plan, which the constraint component of commitment requires the agent to complete. For the second: a person who wants a holiday in Patagonia more than anything else she might do with the money, and knows that she must save for it, but has not decided to go and does not intend to save, is not incoherent; she is merely undecided. Her desire leaves the question open, and nothing in desiring an end requires one to settle on its means.
\end{proof}

\begin{corollary}\label{c:distinct}
Intentions are distinct from any combination of beliefs and desires, and are distinguished by the commitment they involve.
\end{corollary}

\begin{proof}
By Proposition~\ref{p:means} intentions are governed by the means--end requirement and the corresponding beliefs and desires are not. By Principle~\ref{p:func} the attitude types are distinct. The difference consists in the constraint component of commitment, and Propositions~\ref{p:buridan} and~\ref{p:undecided} show that the settling component also distinguishes them.
\end{proof}

\paragraph{Consistency and agglomeration.} The planning conception yields the other rational requirements on intention. Since an intention settles what to do in a world the agent takes herself to inhabit, it must be consistent with what she believes: one cannot intend what one believes impossible, though, as ACT-11 argues, one can intend to try. Since plans are coordinated wholes, the intentions that belong to a single plan must be jointly consistent with belief and are agglomerative: intending $A$ and intending $B$ as parts of one plan commits one to intending $A$ and $B$. Intentions that do not belong to a single plan, such as alternative means pursued in parallel to a single end, are not subject to agglomeration, which is the result ACT-15-a uses to show that an agent may intend each of two outcomes she knows to be incompatible. Neither requirement has a counterpart for desire.

\section{Objections and replies}

\begin{objection}[Cognitivism]
Intention is a belief that one will $A$, formed because one wants it to be true and because the belief, once held, tends to make itself true \citep{Velleman1989,Setiya2007}. The means--end requirement then follows from theoretical rationality: if one believes one will $E$ and that one will $E$ only if one $M$s, consistency of belief requires believing that one will $M$ \citep{Harman1976,Wallace2001}.
\end{objection}

\begin{reply}
The proposal fails in both directions. First, one can intend without believing that one will act. \citet{Bratman1987} describes an agent who intends to stop at the bookstore on the way home but, knowing how absent-minded he is, does not believe that he will; he is not thereby unintending. Second, the derivation of the means--end requirement does not go through. Theoretical consistency forbids believing that one will $E$ while believing that one will not $M$. It does not require forming any belief about $M$ at all, and an agent who has no view about whether she will $M$ is not theoretically inconsistent. But an agent who intends $E$, knows $M$ is necessary, and has not settled whether to $M$ is practically incoherent, as Proposition~\ref{p:means} states. The requirement to settle on means is a requirement on planning, not on belief.
\end{reply}

\begin{objection}[Commitment is a further desire]
What Definition~\ref{d:commit} calls commitment is just a second-order desire: the desire not to reconsider what one has settled, or to be an agent who follows through. Intention is then a predominant first-order desire plus this second-order desire, and reduces after all \citep{Sinhababu2013}.
\end{objection}

\begin{reply}
A desire not to reconsider would be one consideration to be weighed against others in deliberation, and so would be reconsidered whenever deliberation reopened. Commitment does not work that way. It does not add weight to the option intended; it excludes incompatible options from being weighed, in the way that \citet{Raz1975} argued exclusionary reasons do. An agent who intends to take the left branch does not weigh the right branch, finding it outweighed by her desire for consistency; she does not consider it. Moreover, the second-order desire, like any desire, leaves the Buridan case untouched: both walkers may want equally to be agents who follow through.
\end{reply}

\begin{objection}[Bootstrapping]
If intentions constrain reasoning by requiring means, then intending something wrong or foolish generates a reason to take the means to it. That is absurd: one cannot give oneself reasons by deciding.
\end{objection}

\begin{reply}
The means--end requirement is wide in scope: it requires that one not intend $E$, believe $M$ necessary and fail to intend $M$, and it is satisfied as well by abandoning the intention as by taking the means \citep{Broome1999}. It is a requirement of coherence, not a reason to do anything in particular. Commitment so understood creates no reasons; it structures the agent's response to the reasons she has. In this it differs from a promise, which on the view of ETH-16 binds because it creates reasonable reliance in another. An intention binds no one, and its abandonment wrongs no one.
\end{reply}

\begin{objection}[Akrasia and evaluation]
Davidson's later view, that intention is an all-out judgment of desirability, is not a belief--desire view, and the arguments above do not touch it.
\end{objection}

\begin{reply}
The Buridan case does touch it: the walker judges neither branch better, all out, and still intends one. And ACT-04 argues that an agent can intentionally act against her better judgment; on Davidson's view, the akratic agent would have to intend what she judges best while doing something else intentionally, and the intention under which she acts would be left unexplained. On the planning view, the akratic agent forms an intention, or acts on a proximal one, contrary to her judgment, which is what akrasia consists in.
\end{reply}

\section{What the result does not establish}

The argument does not establish that intentions are basic, unanalysable states; it establishes only that they are not reducible to beliefs and desires. It is compatible with an analysis in terms of plans, dispositions and requirements of the kind given in Definition~\ref{d:intention}, and with a functionalist realization of these. Nor does it show that every intention is formed by deliberation: many are formed by picking, as in the Buridan case, and some by habit.

\section{Conclusion}

Two agents with the same beliefs and desires may intend different things, and an agent who wants and expects to do something may not yet intend to do it. Intentions settle practical questions and constrain further reasoning, by excluding incompatible options and requiring means; beliefs and desires do neither. Intentions are plans held with commitment, and they are not reducible to the attitudes of the belief--desire model.

\begin{thebibliography}{99}\small
\bibitem[Audi(1973)]{Audi1973} Audi, R. (1973). Intending. \emph{The Journal of Philosophy}, 70(13), 387--403.
\bibitem[Bratman(1987)]{Bratman1987} Bratman, M.~E. (1987). \emph{Intention, Plans, and Practical Reason}. Harvard University Press.
\bibitem[Broome(1999)]{Broome1999} Broome, J. (1999). Normative requirements. \emph{Ratio}, 12(4), 398--419.
\bibitem[Davidson(1963)]{Davidson1963} Davidson, D. (1963). Actions, reasons, and causes. \emph{The Journal of Philosophy}, 60(23), 685--700.
\bibitem[Davidson(1978)]{Davidson1978} Davidson, D. (1978). Intending. In Y.~Yovel (ed.), \emph{Philosophy of History and Action}. D. Reidel. Reprinted in \emph{Essays on Actions and Events}, Clarendon Press, 1980.
\bibitem[Harman(1976)]{Harman1976} Harman, G. (1976). Practical reasoning. \emph{The Review of Metaphysics}, 29(3), 431--463.
\bibitem[Kavka(1983)]{Kavka1983} Kavka, G.~S. (1983). The toxin puzzle. \emph{Analysis}, 43(1), 33--36.
\bibitem[Raz(1975)]{Raz1975} Raz, J. (1975). \emph{Practical Reason and Norms}. Hutchinson.
\bibitem[Setiya(2007)]{Setiya2007} Setiya, K. (2007). \emph{Reasons without Rationalism}. Princeton University Press.
\bibitem[Sinhababu(2013)]{Sinhababu2013} Sinhababu, N. (2013). The desire-belief account of intention explains everything. \emph{No\^us}, 47(4), 680--696.
\bibitem[Ullmann-Margalit and Morgenbesser(1977)]{UllmannMargalitMorgenbesser1977} Ullmann-Margalit, E. and Morgenbesser, S. (1977). Picking and choosing. \emph{Social Research}, 44(4), 757--785.
\bibitem[Velleman(1989)]{Velleman1989} Velleman, J.~D. (1989). \emph{Practical Reflection}. Princeton University Press.
\bibitem[Wallace(2001)]{Wallace2001} Wallace, R.~J. (2001). Normativity, commitment, and instrumental reason. \emph{Philosophers' Imprint}, 1(3).
\end{thebibliography}

\end{document}
